% !PRINTED 205
\weberpage{205}
%<<CONT
\gesperrt{die Gruppe so transformiren, dass sie die beiden
Substitutionen}
\[
\Theta = \varepsilon x, \quad \psi = \frac{c}{x}
\]
\gesperrt{enthält, worin $\varepsilon$ eine primitive $\nu^{\text{te}}$ Einheitswurzel
bedeutet, und $c$ ein beliebig vorgeschriebener
Werth, z. B. auch $1$, sein kann.}

Die Voraussetzung dieses Satzes ist bei den in \S.\,52 aufgezählten
 Fällen immer für einen der verschiedenen Werthe $\nu$
erfüllt, ausgenommen bei der cyklischen Gruppe und bei der
Diëdergruppe mit $m = 2$.

Endlich beweisen wir noch den folgenden Satz:

2.~\gesperrt{Sind $a, a'$ die Pole einer Substitution $\Theta$ von $G$,
so sind die mit $a$ und $a'$ conjugirten Pole}
\[
b = \chi(a), \quad b' = \chi(a'),
\]
\gesperrt{worin $\chi$ eine beliebige Substitution aus $G$ ist, die
beiden Pole einer und derselben Substitution,
nämlich der Substitution $\chi\Theta\chi^{-1}$.}

Die Richtigkeit ergiebt sich unmittelbar aus dem Anblick
der Gleichungen:
\[
\begin{aligned}
a &= \Theta(a), & \chi\Theta\chi^{-1}\chi(a) &= \chi\Theta(a) = \chi(a) \\
a' &= \Theta(a'), & \chi\Theta\chi^{-1}\chi(a') &= \chi\Theta(a') = \chi(a').
\end{aligned}
\]

\begin{center}
\S.\,54.

\gesperrt{Die Grundformen.}
\end{center}

Um zu der endgültigen Aufstellung aller endlichen Gruppen
$G$ zu gelangen, ist es nothwendig, auf die Invarianten der entsprechenden
 binären Substitutionsgruppen näher einzugehen.

Wir müssen daher neben den linearen gebrochenen Substitutionen


\weq{1}{y = \frac{\alpha x + \beta}{\gamma x + \delta},}

die wir im \S.\,49 mit $A'''$ bezeichnet haben, noch die binären
Substitutionen

\weq{2}{(y_1, y_2) = \begin{pmatrix}\alpha, & \beta \\ \gamma, & \delta\end{pmatrix}(x_1, x_2)}
%CONT>>
